\section{引言}
\label{sec:intro}

\begin{figure}[t]
\centering
\includegraphics[width=\linewidth]{figures/fig1.pdf}
\caption{TwinRouterBench 概览。该基准提供用于离线路由器开发的快速静态赛道，以及用于端到端验证的在线动态赛道。静态赛道覆盖五类工作负载中 520 个实例的 970 条步级记录，每条都带有经执行验证的目标层级；动态赛道在 SWE-bench Verified 上运行路由器，并记录实际 API 成本。}
\label{fig:overview}
\end{figure}

LLM 路由器决定应由哪个模型处理每一次调用~\citep{feng2024graphrouter,stripelis2024tensoropera}。在生产环境的多轮智能体中，每次调用都带有丰富的前缀---聊天历史、检索段落、shell 日志、工具输出和部分代码修改---而在第 $i$ 步作出的低价选择，可能只会在后续步骤失败时才暴露问题。现有基准在彼此独立的一次性提示词上评测路由；它们既不暴露这些前缀，也不检验某个逐步决策是否会破坏之后的步骤。

现代应用中昂贵的单元很少是一次性查询。编码智能体通过多轮探索、修改和验证来解决 GitHub issue~\citep{jimenez2023swebench}；工具系统在串行和并行调用之间维护状态~\citep{patil2025bfcl}；检索和摘要工作负载携带长的多轮上下文~\citep{katsis2025mtrag,zhong2021qmsum}。较便宜的模型或许足以完成文件查找，而后续编写补丁的步骤可能需要最强层级。一个有用的步级基准必须覆盖这些工作负载，以执行结果而非仅依赖模型判断来确定标签，并同时支持低成本、确定性的离线评分与以实际 API 成本为最终验证的端到端执行。因此，我们维护两个彼此区分的赛道。

RouterBench、LLMRouterBench 和 RouterArena 在完整的一次性提示词上评测查询级路由~\citep{hu2024routerbench,li2026llmrouterbench,lu2025routerarena}。此类路由器仍可通过把完整对话历史作为单个查询输入而应用在轨迹中间，但其训练与评测数据均为一次性提示词，所以无法说明在第 $i$ 步使用低价模型会不会破坏后续步骤。我们将每次 LLM 调用前路由器观察到的完整前缀---聊天历史、检索段落、工具输出、shell 轨迹和代码 diff---作为一等输入，并通过将任务运行至结束来验证每个标签：若把某一步降级后，轨迹仍以相同步数通过，我们便推断被替换的步骤已被正确处理。TRIM 研究了数学问题中推理步骤的升级~\citep{kapoor2026trim}；我们的场景则是多轮智能体。

构建这样的基准面临两项主要挑战。第一，将前沿模型直接用作路由器并不可靠：在我们的 Claude Opus 4.6 诊断中，单模板 LLM 路由器在 147 个经验证的 \texttt{high} SWE-bench 步骤中只预测了 7 个 \texttt{high}，并使全部 SWE 轨迹失败。第二，多轮轨迹带来组合式标签复杂度：对于 $n$ 个候选模型和 $m$ 个步骤，即使轨迹长度固定，穷举验证也需要 $n^{m}$ 次运行，因为第 $i{+}1$ 步的前缀取决于第 $i$ 步选择的模型。因此，实用的标签构建必须以可负担的协议近似这一空间。

我们提出 \textbf{TwinRouterBench}，一个用于离线开发的快速静态赛道和一个用于端到端验证的在线动态赛道组成的步级路由基准。

\paragraph{贡献。}
\begin{itemize}[leftmargin=*, itemsep=1pt, topsep=2pt]
\item \textbf{面向真实路由监督的快速静态赛道。}静态赛道覆盖最能从步级路由受益的工作负载（长上下文、多轮、工具使用、RAG、摘要和代码修复），并暴露每次 LLM 调用前路由器所见的精确前缀。其 970 个标签是在固定的降级--级联协议下得到的执行验证估计，包含层级池级联验证、混合前缀重建，以及针对多步可执行工作负载的人工审计。评分只对标签、轨迹归属和 token 成本做确定性算术计算：评测耗时仅毫秒级，且不需要在线评测端 LLM 裁判。这些数据也支持训练；在我们评测的策略中，基于静态标签训练的逻辑回归路由器实现了最佳的成本--成功权衡。
\item \textbf{面向软件智能体路由的在线动态赛道。}动态 harness 支持完整的 500 个 SWE-bench Verified 样例；本文报告与静态 SWE 监督集不重叠的 100 个留出样例评测。每次 LLM 请求中，路由器从锁定模型池选择一个具体模型；评分使用官方 SWE-bench 解决判定、实际提供商支出、缓存计费，并在排行榜账单中加入固定的未解决实例惩罚（\S\ref{sec:dynamic-scoring}），且没有评测端 LLM 裁判。
\item \textbf{双赛道开发与验证闭环。}静态赛道提供低成本、可训练的监督以支持快速迭代；动态赛道检验所得策略能否经受真实执行。我们并列报告两者，使离线证据与端到端证据保持可区分。
\end{itemize}
